% -----------------------------------------------------------------------------
% Document class
% scrartcl = clean article class for reports; bibliography in TOC, good tables
% -----------------------------------------------------------------------------
\documentclass[
  bibliography=totoc,
  captions=tableheading,
]{scrartcl}

% Fix minor KOMA compatibility issues
\usepackage{scrhack}

% Warn if recompilation needed (refs, TOC, etc.)
\usepackage[aux]{rerunfilecheck}

% Core math packages (align, symbols, extensions)
\usepackage{amsmath}
\usepackage{amssymb}
\usepackage{mathtools}

% Modern font handling (requires lualatex/xelatex)
\usepackage{fontspec}
\recalctypearea{} % recompute layout after font setup

% Language settings (hyphenation, captions)
\usepackage[english]{babel}

% Modern math fonts + ISO style (clean scientific notation)
\usepackage[
  math-style=ISO,
  bold-style=ISO,
  sans-style=italic,
  nabla=upright,
  partial=upright,
  mathrm=sym,
]{unicode-math}

\setmathfont{Latin Modern Math} % default math font
\setmathfont{XITS Math}[range={scr, bfscr}] % nicer script letters
\setmathfont{XITS Math}[range={cal, bfcal}, StylisticSet=1]

% Numbers & units formatting (\num, \SI, etc.)
\usepackage[
  locale=US,
  separate-uncertainty=true,
  per-mode=symbol-or-fraction,
  detect-all,
]{siunitx}

% Quotes (recommended with biblatex)
\usepackage[english=american]{csquotes}

% Float placement control (figures/tables positioning)
\usepackage{float}
\floatplacement{figure}{htbp}
\floatplacement{table}{htbp}

% Keep figures inside sections (prevents drifting)
\usepackage[section,below]{placeins}

% text wrapping around figures
\usepackage{wrapfig}      

% Caption styling (bold labels, smaller font)
\usepackage[
  labelfont=bf,
  font=small,
  width=0.9\textwidth,
]{caption}

% Subfigures (compare plots side-by-side)
\usepackage{subcaption}

% Include images
\usepackage{graphicx}

% Better tables (clean formatting + number alignment)
\usepackage{tabularray}
\UseTblrLibrary{booktabs, siunitx}

% Typography improvements (spacing, kerning)
\usepackage{microtype}

% Bibliography system (modern replacement for BibTeX)
\usepackage[backend=biber]{biblatex}
\addbibresource{../../../latex/references.bib}

% Clickable links + PDF metadata
\usepackage[
  unicode,
  pdfusetitle,
  pdfcreator={},
  pdfproducer={},
]{hyperref}

% Better PDF bookmarks
\usepackage{bookmark}

% Custom commands (needed for definitions below)
\usepackage{expl3}
\usepackage{xparse}

% Upright e (useful for exp, softmax, etc.)
\ExplSyntaxOn
\NewDocumentCommand \E {} {\symup{e}}
\ExplSyntaxOff

% Upright differential d (clean integrals)
\ExplSyntaxOn
\NewDocumentCommand \dif {m} {\mathinner{\symup{d} #1}}
\ExplSyntaxOff

% Vector command (bold vectors in math mode)
\ExplSyntaxOn
\let\vaccent=\v
\RenewDocumentCommand \v {}{
  \TextOrMath{\vaccent}{\symbf}
}
\ExplSyntaxOff

% Better spacing for scalable brackets
\usepackage{mleftright}

% Shortcuts for scalable brackets: \l( ... \r)
\ExplSyntaxOn
\let\ltext=\l
\RenewDocumentCommand \l {}{\TextOrMath{\ltext}{\mleft}}
\let\raccent=\r
\RenewDocumentCommand \r {}{\TextOrMath{\raccent}{\mright}}
\ExplSyntaxOff

% Absolute value and norm (auto-scaling with *)
\DeclarePairedDelimiter{\abs}{\lvert}{\rvert}
\DeclarePairedDelimiter{\norm}{\lVert}{\rVert}

% Definition equality symbols
\newcommand{\defeq}{\vcentcolon=}
\newcommand{\eqdef}{=\vcentcolon}

% No paragraph indent (clean report style)
\setlength{\parindent}{0pt}

\title{Task 06: Minimal 1D Diffusion Model}
\author{Akram Aki}
\date{\today}

\begin{document}

\maketitle

\section{Introduction}

The goal of this task was to implement a minimal diffusion model for one-dimensional data.
The target distribution is a mixture of two Gaussian modes centered at $-4$ and $4$ with unit standard deviation.
This setting is intentionally simple, but it makes the main idea of diffusion models visible: a fixed forward process gradually adds noise, and a learned reverse process transforms noise back into samples from the target distribution.

\section{Method}

The forward diffusion process used a fixed variance schedule with
\begin{equation*}
    \beta_t = 0.02, \qquad T = 250.
\end{equation*}
For a clean sample $x_0$, noisy samples were generated directly using
\begin{equation*}
    x_t = \sqrt{\bar{\alpha}_t} x_0 + \sqrt{1 - \bar{\alpha}_t}\epsilon,
    \qquad \epsilon \sim \mathcal{N}(0, 1),
\end{equation*}
where $\bar{\alpha}_t = \prod_{s=1}^{t}(1-\beta_s)$.

The denoising model was a small fully connected MLP. It receives the noisy scalar value $x_t$ and the normalized timestep $t/T$, and predicts the noise $\epsilon$ that was added.
The model was trained with mean squared error loss and Adam.
After training, samples were generated by starting from standard Gaussian noise and applying the learned reverse denoising process step by step.
The notebook was run for two target mixtures: a balanced $50/50$ mixture and an imbalanced $25/75$ mixture.
The implementation is available in the accompanying \href{https://github.com/AkramAki/Advanced_DeepLearning_Seminar/tree/main/tasks/task_06/code/task_06.ipynb}{GitHub repository notebook}.

\newpage
\section{Results}

\autoref{fig:forward-process} shows the fixed forward diffusion process for the imbalanced target distribution, where the left mode has probability $0.25$ and the right mode probability $0.75$.
The two modes are clearly visible at early timesteps. As the timestep increases, Gaussian noise dominates and the original mixture structure gradually disappears.
This confirms that the forward process performs the intended corruption from data to noise.

\begin{figure}[htbp]
    \centering
    \includegraphics[width=0.9\textwidth]{../build/forward_diffusion_steps.pdf}
    \caption{Forward diffusion snapshots for the imbalanced mixture. The original two-mode structure is gradually destroyed as the timestep increases.}
    \label{fig:forward-process}
\end{figure}

\autoref{fig:reverse-process} shows the learned reverse diffusion process for the same imbalanced target distribution.
The initial samples are close to standard Gaussian noise. During reverse diffusion, the distribution gradually separates into two modes, and the larger right mode is preserved.

\begin{figure}[htbp]
    \centering
    \includegraphics[width=0.9\textwidth]{../build/reverse_diffusion_steps.pdf}
    \caption{Reverse diffusion snapshots for the imbalanced mixture with $25\%$ of samples in the left mode and $75\%$ in the right mode.}
    \label{fig:reverse-process}
\end{figure}

\autoref{fig:final-comparison} compares the final generated distribution with the true target distribution for the two notebook runs.
In the balanced case, the model recovers two modes with similar probability mass.
In the imbalanced case, it also learns the unequal mode weights, showing that the same simple diffusion setup is not restricted to symmetric targets.

\begin{figure}[htbp]
    \centering
    \begin{subfigure}{0.48\textwidth}
        \centering
        \includegraphics[width=\textwidth]{../build/final_generated_vs_target_50.pdf}
        \caption{Balanced mixture, $50/50$.}
        \label{fig:final-balanced}
    \end{subfigure}
    \hfill
    \begin{subfigure}{0.48\textwidth}
        \centering
        \includegraphics[width=\textwidth]{../build/final_generated_vs_target.pdf}
        \caption{Imbalanced mixture, $25/75$.}
        \label{fig:final-imbalanced}
    \end{subfigure}
    \caption{Final generated samples compared with the true target distribution. The model captures the two-mode structure in both settings and also follows the unequal mode probabilities in the imbalanced case.}
    \label{fig:final-comparison}
\end{figure}

\section{Conclusion}

The implemented 1D diffusion model successfully learned to generate samples from a two-Gaussian target distribution.
Although the model is very small and uses only scalar timestep conditioning, it captures both the balanced and imbalanced mixtures reasonably well.
The main limitation is that the setup is intentionally minimal; possible improvements include richer timestep embeddings, a learned or scheduled $\beta_t$, longer training, and quantitative distribution metrics in addition to visual comparisons.

\printbibliography

\end{document}
